% TOP_PROBLEM: {"id": 3167, "title": "r-regular graphs are not uniquely hamiltonian.", "category_id": 3, "category": {"id": 3, "name": "graph_theory", "display_name": "Graph Theory", "description": "Problems involving graphs, networks, and their properties.", "slug": "graph-theory", "order_index": 3, "created_at": "2026-07-31T15:26:25.671Z"}, "status": "open", "problem_number": "OPG-480", "set_id": 12, "statusQualification": "The graph is explicitly simple, as required by the source comments excluding multigraph counterexamples."}
% FIRST_POSTED: 2026-10-01
% ATTEMPT_STATUS: unresolved
% UnsolvedMath ID 3167; source catalogue code OPG-480
% !TeX program = lualatex
% GPT-6 Astra Ultra research attempt; independent partial work, 2026-10-01.
\documentclass[11pt,a4paper]{article}
\usepackage[margin=25mm]{geometry}
\usepackage{fontspec}
\setmainfont{lmroman10-regular.otf}[BoldFont=lmroman10-bold.otf,
 ItalicFont=lmroman10-italic.otf,BoldItalicFont=lmroman10-bolditalic.otf]
\usepackage{amsmath,amssymb,mathtools}
\usepackage{unicode-math}
\setmathfont{latinmodern-math.otf}
\AtBeginDocument{\let\setminus\smallsetminus}
\usepackage{xurl,hyperref}
\hypersetup{colorlinks=true,urlcolor=blue}
\setcounter{secnumdepth}{0}
\setlength{\parindent}{0pt}
\setlength{\parskip}{5pt}
\setlength{\emergencystretch}{3em}
\begin{document}
\section*{r-regular graphs are not uniquely hamiltonian.}\label{3167}
{\small UnsolvedMath ID 3167 · OPG-480 · Graph Theory · OpenGarden}
\subsection{Definitions and mathematical statement}
Let $r\geq3$ be an integer and let $G=(V,E)$ be a finite simple
undirected $r$-regular graph: every vertex has exactly $r$
neighbours. A Hamilton cycle is a connected spanning
$2$-regular subgraph, equivalently a cycle passing through
all vertices once. Count Hamilton cycles by their edge sets.
Changing the starting vertex or reversing the direction of
traversal does not create a new cycle.

The graph is uniquely Hamiltonian when it has exactly one
Hamilton cycle. The conjecture asserts that no finite simple
$r$-regular graph is uniquely Hamiltonian for any $r>2$.
Writing $h(G)$ for the number of Hamilton cycles, its target
is
\[
 h(G)\neq1\qquad
 \text{for every finite simple regular graph of degree }r\geq3.
\]
Equivalently, whenever such a $G$ has a Hamilton cycle $C$,
there must be another Hamilton cycle $C'$ with
$E(C')\neq E(C)$.

The assertion permits graphs with no Hamilton cycle; it
does not assert that regularity alone guarantees one. The
two cycles may share vertices and edges and need not be
edge-disjoint. No connectivity assumption beyond that
forced by existence of the first Hamilton cycle is needed.
Simplicity is essential to the intended source question:
parallel edges are excluded, as emphasized in its discussion.
Degree two is omitted because a connected simple cycle is
itself a uniquely Hamiltonian $2$-regular graph.
\subsection{Short English statement}
Must every Hamiltonian regular simple graph of degree greater than two have at least two distinct Hamilton cycles?
\subsection{Research attempt: parity for odd degrees and a reduction to degree four}
\paragraph{Scope and literature check.}
GPT-6 Astra Ultra, 1 October 2026. Hamilton cycles are counted as
undirected edge sets. The task is to exclude exactly one, not to
prove that a regular graph is Hamiltonian. Thomason's parity theorem
[S3] settles the odd-degree case. We give the path-rotation argument
explicitly and then prove that a resolution for simple four-regular
graphs would settle every remaining even degree. These are established
reductions recovered here, not new solutions.

\paragraph{1. A finite auxiliary graph of Hamilton paths.}
Suppose every vertex of $G$ has odd degree, and fix an edge $ab$.
Form an auxiliary graph whose vertices are all ordered Hamilton paths
\[
 P=(v_1,v_2,\ldots,v_n),\qquad v_1=a,\ v_2=b.
\]
If the final vertex $z=v_n$ is adjacent to $v_i$ for
$2\le i\le n-2$, rotate the final segment to obtain
\[
 P'=(v_1,\ldots,v_i,v_n,v_{n-1},\ldots,v_{i+1}).
\]
This remains a Hamilton path with the fixed initial edge. The move
is reversible, using the removed edge $v_iv_{i+1}$, and distinct
choices of $i$ give distinct neighbours. There are no loops. Its
degree at $P$ is therefore
\[
 d_{\rm aux}(P)=\deg_G(z)-1-\mathbf1_{\{za\in E(G)\}}.
\]
The subtracted one accounts for the predecessor of $z$; adjacency
to $a$ cannot be used because rotations must preserve the initial
edge. Every other neighbour is eligible.

Since $\deg_G(z)$ is odd, the auxiliary degree is odd precisely
when $za$ is an edge. The handshaking lemma says that the number
of odd-degree vertices of this finite auxiliary graph is even.
Closing such a path with $za$ gives a Hamilton cycle containing $ab$.
Conversely, orienting any Hamilton cycle containing $ab$ to begin
$a,b$ and removing its final edge to $a$ gives exactly one such
ordered path. Hence the number of Hamilton cycles containing $ab$
is even. If a Hamilton cycle exists, choose $ab$ on it; the count
is positive and even, so a second edge-distinct cycle exists. This
proves the conjecture for every odd regular degree, and indeed for
all graphs with every degree odd.

\paragraph{2. Reducing all even degrees to four.}
Let $G$ be simple, $r$-regular with even $r\ge4$, and let $C$ be
a Hamilton cycle. Its edge complement $R=G-E(C)$ is $(r-2)$-regular.
We show it has a spanning two-factor. Orient an Euler tour of each
connected component of $R$. Every vertex then has indegree and
outdegree $d=(r-2)/2\ge1$.

Make a bipartite graph with a left and right copy of every vertex,
and an edge from the left copy of $u$ to the right copy of $v$
for each directed edge $u\to v$. This graph is $d$-regular on both
sides. Counting edges from any left subset proves Hall's inequality,
so it has a perfect matching. The selected directed edges give one
incoming and one outgoing edge at each original vertex, hence a
disjoint union of directed cycles spanning $V(G)$. No loop exists,
and a directed two-cycle would require both orientations of one
original simple edge, which were not included. Forgetting orientations
therefore gives an ordinary spanning two-factor $F\subseteq R$.

The union $C\cup F$ is a simple four-regular spanning subgraph of
$G$. A second Hamilton cycle there would also be a second Hamilton
cycle in $G$. Thus the four-regular case is sufficient for every
even regular degree. Together with the parity proof, it suffices for
the entire canonical question.

\paragraph{Verification and the exact parity barrier.}
For even original degrees, the displayed auxiliary-degree formula
makes closing paths have even auxiliary degree and nonclosing paths
odd. The handshaking lemma then controls the wrong class of paths;
it no longer forces a second Hamilton cycle. This is a precise failure
of extending the same parity argument, not a counterexample. The
script [S4] checks Hamilton-cycle counts through one fixed edge of
$K_4$ and $K_6$, obtaining two and twenty-four respectively. Those
small counts illustrate the odd-degree theorem but do not address
the missing four-regular case.

\subsection{Final status}
UNRESOLVED. The odd-degree parity theorem is proved and the remaining
regular cases are reduced to degree four. The four-regular case is
not resolved. Completion estimate: no defensible global percentage.

\subsection{Sources}
\begin{itemize}
\item[S1] ULAM.AI and the UnsolvedMath Contributors, supplied problems.json, record 3167 (OPG-480), OpenGarden collection. Catalogue identity and source transcription; CC BY 4.0.
\url{https://huggingface.co/datasets/ulamai/UnsolvedMath}.
\item[S2] Open Problem Garden, r-regular graphs are not uniquely hamiltonian.. Original problem and discussion; the mathematical formulation below is expanded from the supplied source record.
\url{http://www.openproblemgarden.org/op/uniquely_hamiltonian_graphs}.
\item[S3] A. G. Thomason, \emph{Hamiltonian Cycles and Uniquely Edge Colourable Graphs}, Annals of Discrete Mathematics 3 (1978), 259--268. Author-hosted paper checked 1 October 2026.
\url{https://www.dpmms.cam.ac.uk/~agt2/HC.pdf}.
\item[S4] GPT-6 Astra Ultra, finite Hamilton-cycle count checks, 1 October 2026: \texttt{scripts/check\_attempt\_batch02\_connectivity\_2026\_10\_01.py} in this repository; run with Python 3.
\end{itemize}
\end{document}
